\begin{pastexam}[label={px:ncu-109-dsa-12}]{中央資工所 109 學年度 申論題第 1 題}
（申論）

(10\%) Suppose that two stacks, stack 0 and stack 1, are both mapped to array mem[MEMORY\_SIZE].
Algorithm PUSH adds an item into stack i, where 0$\le$i$<$ 2, as long as the total number of items
in these two stacks is less than MEMORY\_SIZE$-$ 1. Please complete algorithm PUSH.
\begin{minted}[fontsize=\small]{c}
int mem[MEMORY_SIZE];   /*global variable */
int top[2];    /*global variable */
void PUSH (integer item, integer i)
{

}
\end{minted}
\end{pastexam}

\begin{pastsoln}
題目有幾個地方要先講清楚。
\begin{itemize}
\item 題目沒有給 \texttt{top} 的初值。 照兩個堆疊共用一個陣列的標準做法（資結書第~6.5~節）：
堆疊 $0$ 從 \texttt{mem[0]} 往右長，堆疊 $1$ 從 \texttt{mem[MEMORY\_SIZE-1]} 往左長，\texttt{top[i]}
指向堆疊 $i$ 的頂端元素。 空的時候 \texttt{top[0] = -1}、\texttt{top[1] = MEMORY\_SIZE}，由程式
開始時的初始化設定。 這時堆疊 $0$ 有 \texttt{top[0] + 1} 個元素，堆疊 $1$ 有
\texttt{MEMORY\_SIZE - top[1]} 個元素。
\item 「total number of items is less than MEMORY\_SIZE $-$ 1」照字面是：放進去\textbf{之前}，兩個堆疊
的總數要小於 \texttt{MEMORY\_SIZE - 1}，才可以放。 這樣最多只放得進 \texttt{MEMORY\_SIZE - 1} 個，
陣列會留一格不用。 下面照字面寫，最後再說只要改哪一行就能用滿整個陣列（一般課本的版本）。
\item 原題的參數型別寫成 \texttt{integer}，不是 C 的型別，下面改成 \texttt{int}；滿了的時候印出訊息後
直接返回。
\end{itemize}

\begin{minted}[fontsize=\small]{c}
void PUSH(int item, int i)
{
    int total = (top[0] + 1) + (MEMORY_SIZE - top[1]);  /* 兩個堆疊目前共有幾個 */
    if (total >= MEMORY_SIZE - 1) {   /* 題目：total < MEMORY_SIZE - 1 才能放 */
        fprintf(stderr, "stack full\n");
        return;
    }
    if (i == 0)
        mem[++top[0]] = item;     /* 堆疊 0 往右長：先加一再放 */
    else
        mem[--top[1]] = item;     /* 堆疊 1 往左長：先減一再放 */
}
\end{minted}

\textbf{正確性}。 任何時候堆疊 $0$ 佔 \texttt{mem[0..top[0]]}，堆疊 $1$ 佔
\texttt{mem[top[1]..MEMORY\_SIZE-1]}，兩段不重疊，中間 \texttt{mem[top[0]+1..top[1]-1]} 是空的，
共 \texttt{top[1] - top[0] - 1} $=$ \texttt{MEMORY\_SIZE} $-$ \textit{total} 格。 通過檢查時
$\mathit{total} \le \texttt{MEMORY\_SIZE} - 2$，中間至少空兩格，所以不論放進哪一個堆疊，新的位置
\texttt{top[0]+1} 或 \texttt{top[1]-1} 都落在空的那一段裡，不會蓋掉另一個堆疊的資料，也不會超出
陣列。 放完之後兩段仍然不重疊，性質保持。 \texttt{++top[0]} 先加一再放、\texttt{--top[1]} 先減一再放，
讓 \texttt{top[i]} 一直指向頂端元素。

\textbf{時間}。 只有常數個運算，$O(1)$。

\textbf{用滿整個陣列的版本}。 若題意其實是「陣列還有空位就能放」（兩個頂端碰在一起才算滿，資結書
第~6.5~節的寫法），把檢查改成
\begin{minted}[fontsize=\small]{c}
    if (top[0] + 1 == top[1]) {     /* 中間一格都不剩，total == MEMORY_SIZE */
\end{minted}
其他不變。 兩種寫法差別只在最多能放 \texttt{MEMORY\_SIZE - 1} 個還是 \texttt{MEMORY\_SIZE} 個；作答時
把採用的條件寫清楚即可。 這種從兩端往中間長的配置，好處是只要總數還沒到上限，任何一個堆疊都還能
放，不會發生一個堆疊滿了、另一個還很空的情形。
\end{pastsoln}
